\section{Verification Trigger Frequency and Path Generation Statistics}
\label{app:trigger_stats}

\subsection{Verification Trigger Frequency}

Table~\ref{tab:trigger_rate} reports the average verification trigger rate across benchmarks and backbone models. A question is counted as ``triggered'' if at least one verification call is made during its debate.

\begin{table}[h]
\centering
\resizebox{\columnwidth}{!}{%
\begin{tabular}{lccc}
\toprule
\textbf{Benchmark} & \textbf{GPT-4o-mini} & \textbf{Qwen3-8B} & \textbf{Llama-3.3-70B} \\
\midrule
GSM8K       & 12.2\% & 12.1\% &  5.0\% \\
MMLU        & 25.5\% & 31.8\% & 17.2\% \\
MATH500     & 40.4\% & 29.3\% & 30.4\% \\
GPQA        & 53.7\% & 58.1\% & 35.0\% \\
AIME 2024   & 91.1\% & 84.4\% & 78.9\% \\
AIME 2025   & 92.2\% & 82.2\% & 88.9\% \\
\bottomrule
\end{tabular}
}
\caption{Average verification trigger rate (\%) per benchmark and backbone model. Benchmarks are ordered by increasing difficulty.}
\label{tab:trigger_rate}
\end{table}

The trigger rates exhibit a clear correlation with task difficulty: the verification agent fires on fewer than 13\% of GSM8K problems but on over 78\% of AIME problems across all backbones. This pattern confirms that the verification module is adaptive rather than indiscriminate---it allocates additional computational effort selectively to problems where agent disagreement genuinely arises. As a consequence, the token overhead reported in Appendix~\ref{app:cost} is concentrated on the hardest tasks, precisely where the accuracy gains from verification are largest.

\subsection{Path Generation Statistics}

Table~\ref{tab:path_gen} reports the average number of distinct reasoning paths generated per question by the Path Generation module ($\Phi_{\text{gen}}$). The target allocation is three paths (one per agent).

\begin{table}[h]
\centering
\resizebox{\columnwidth}{!}{%
\begin{tabular}{lccc}
\toprule
\textbf{Benchmark} & \textbf{GPT-4o-mini} & \textbf{Qwen3-8B} & \textbf{Llama-3.3-70B} \\
\midrule
GSM8K     & 3.00 & 2.98 & 2.84 \\
MMLU      & 3.00 & 3.00 & 3.00 \\
MATH500   & 2.95 & 2.92 & 2.33 \\
GPQA      & 3.00 & 3.00 & 3.00 \\
AIME 2024 & 2.91 & 3.00 & 2.69 \\
AIME 2025 & 3.00 & 3.00 & 2.62 \\
\bottomrule
\end{tabular}
}
\caption{Average number of reasoning paths generated per question. The target is 3 paths (one per agent).}
\label{tab:path_gen}
\end{table}

Path generation is highly reliable: GPT-4o-mini and Qwen3-8B consistently produce the full complement of three paths across all benchmarks. Llama-3.3-70B occasionally generates fewer paths on mathematical tasks (e.g., 2.33 on MATH500), reflecting variation in instruction-following capacity across backbone models; the framework gracefully falls back to the available paths in such cases without disrupting the debate pipeline.

\subsection{Path Independence: Single-path Rate vs.\ Path Coverage Rate}

To directly assess whether the generated paths are genuinely complementary rather than redundant paraphrases, we compare two metrics computed from each agent's \emph{pre-debate} initial answer (i.e., before any inter-agent interaction):

\begin{itemize}
    \item \textbf{Single-path Rate}: fraction of correct initial answers averaged over all paths and questions, equivalent to the accuracy of a single randomly selected path.
    \item \textbf{Path Coverage Rate}: fraction of questions for which \emph{at least one} of the three paths yields a correct initial answer.
\end{itemize}

If the paths were near-identical, coverage would barely exceed the single-path rate. A large gap instead indicates that different paths succeed on different questions---the defining signature of genuine complementarity.

\begin{table}[h]
\centering
\resizebox{\columnwidth}{!}{%
\begin{tabular}{llcc}
\toprule
\textbf{Benchmark} & \textbf{Model} & \textbf{Single-path} & \textbf{Coverage} \\
\midrule
{GSM8K}
 & GPT-4o-mini    & 91.1\% & 95.3\% \\
 & Qwen3-8B       & 90.8\% & 95.2\% \\
 & Llama-3.3-70B  & 95.0\% & 96.0\% \\
\midrule
{MMLU}
 & GPT-4o-mini    & 79.4\% & 89.3\% \\
 & Qwen3-8B       & 76.9\% & 88.2\% \\
 & Llama-3.3-70B  & 84.8\% & 91.9\% \\
\midrule
{MATH500}
 & GPT-4o-mini    & 66.7\% & 77.4\% \\
 & Qwen3-8B       & 75.8\% & 84.6\% \\
 & Llama-3.3-70B  & 69.3\% & 76.4\% \\
\midrule
{GPQA}
 & GPT-4o-mini    & 43.6\% & 64.0\% \\
 & Qwen3-8B       & 43.5\% & 64.0\% \\
 & Llama-3.3-70B  & 58.4\% & 74.8\% \\
\midrule
{AIME 2024}
 & GPT-4o-mini    &  6.7\% & 11.1\% \\
 & Qwen3-8B       & 13.7\% & 21.1\% \\
 & Llama-3.3-70B  & 14.1\% & 20.0\% \\
\midrule
{AIME 2025}
 & GPT-4o-mini    &  8.1\% & 13.3\% \\
 & Qwen3-8B       & 20.7\% & 32.2\% \\
 & Llama-3.3-70B  &  5.2\% & 12.2\% \\
\bottomrule
\end{tabular}
}
\caption{Single-path Rate vs.\ Path Coverage Rate (\%) measured from pre-debate initial answers. Benchmarks are ordered by increasing difficulty.}
\label{tab:path_coverage}
\end{table}

As shown in Table~\ref{tab:path_coverage}, Path Coverage Rate consistently and substantially exceeds Single-path Rate across all benchmarks and backbones. The gap is most pronounced on GPQA ($+$16--20 percentage points), where expert-domain problems admit only a narrow set of valid reasoning approaches and paths can succeed or fail independently. Moderate but consistent gaps also appear on MATH500 ($+$7--9 pp) and MMLU ($+$7--11 pp). Even on the hardest AIME tasks, where absolute accuracy is low, coverage exceeds single-path rate by 5--12 pp, confirming that the three paths explore genuinely distinct solution trajectories rather than converging on the same (often incorrect) approach. The relatively small gap on GSM8K ($+$1--4 pp) is expected: when problems are sufficiently easy, most paths independently reach the correct answer, leaving little room for complementary coverage. Together, these results confirm that $\Phi_{\text{gen}}$ produces paths that are functionally independent, providing the debate with a diverse pool of reasoning candidates from which the best can be selected and refined.
